% TOP_PROBLEM: {"id": 2667, "title": "Kirby Problem 1.8", "category_id": 7, "category": {"id": 7, "name": "topology", "display_name": "Topology", "description": "Properties preserved under continuous deformations.", "slug": "topology", "order_index": 7, "created_at": "2026-07-31T15:26:25.671Z"}, "status": "open", "problem_number": "KP-1.8", "set_id": 11}
% FIRST_POSTED: 2026-10-01
% ENTRY_KIND: statement_only
% UnsolvedMath ID 2667; source catalogue code KP-1.8
% !TeX program = lualatex
% Definition and sources only; this file is not an LLM solution attempt.
\documentclass[11pt,a4paper]{article}
\usepackage[margin=25mm]{geometry}
\usepackage{fontspec}
\setmainfont{lmroman10-regular.otf}[BoldFont=lmroman10-bold.otf,
 ItalicFont=lmroman10-italic.otf,BoldItalicFont=lmroman10-bolditalic.otf]
\usepackage{amsmath,amssymb,mathtools}
\usepackage{unicode-math}
\setmathfont{latinmodern-math.otf}
\AtBeginDocument{\let\setminus\smallsetminus}
\usepackage{xurl,hyperref}
\hypersetup{colorlinks=true,urlcolor=blue}
\setcounter{secnumdepth}{0}
\setlength{\parindent}{0pt}
\setlength{\parskip}{5pt}
\setlength{\emergencystretch}{3em}
\begin{document}
\section*{Kirby Problem 1.8}\label{2667}
{\small UnsolvedMath ID 2667 · KP-1.8 · Topology · Kirby's Problems in Low-Dimensional Topology}
\subsection{Definitions and mathematical statement}
Let $Y$ be an oriented integral homology $3$-sphere, meaning
$H_*(Y;\mathbb Z)\cong H_*(S^3;\mathbb Z)$, and let
$L=L_1\cup\cdots\cup L_m$ be an ordered oriented link in $Y$.
Put $X=Y\setminus\operatorname{int}\nu L$.
Sending the $m$ oriented meridians to the standard basis gives
$\pi_1(X)\twoheadrightarrow\mathbb Z^m$ and a corresponding regular cover
$\widetilde X$. Its Alexander module $H_1(\widetilde X;\mathbb Z)$ is a module over
$R_m=\mathbb Z[t_1^{\pm1},\ldots,t_m^{\pm1}]$ by deck transformations.

The multivariable Alexander polynomial $\Delta_L$ is the order of this module:
the greatest common divisor of the maximal minors in a finite presentation,
with value zero if the module has positive rank. It is defined up to a unit
$\pm t_1^{a_1}\cdots t_m^{a_m}$; the zero module has order one.

For each $m$, characterize precisely the Laurent polynomials that occur as
$\Delta_L$ for links in $S^3$. Ask the same realization question when integral
homology spheres are allowed as ambient manifolds. A characterization must give
both necessary conditions and realizations for every polynomial satisfying them,
respecting the stated unit ambiguity.
\subsection{Short English statement}
Which multivariable Laurent polynomials are Alexander polynomials of links in a sphere or a homology sphere?
\subsection{Sources}
\begin{itemize}
\item[S1] ULAM.AI and the UnsolvedMath Contributors, supplied problems.json, record 2667 (KP-1.8), Kirby's Problems in Low-Dimensional Topology. Catalogue identity and source transcription; CC BY 4.0.
\url{https://huggingface.co/datasets/ulamai/UnsolvedMath}.
\item[S2] K3: A New Problem List in Low-Dimensional Topology, Problem 1.8. Expanded formulation of the supplied catalogue statement.
\url{https://math.berkeley.edu/sites/default/files/surv-295-ruberman-watermarked-author-pdf.pdf}.
\end{itemize}
\end{document}
